\section{The PIOP: collapse and shape the constraint system}\label{sec:piop}

This section is the zero-check.  Its input is the integer relation from
\cref{sec:arith}; its output is one linear claim about the virtual bit witness.
For SHA-256 the relation is linear, so the nonlinear outer sumcheck disappears.
The inner sumcheck remains: the public row collapse produces an arbitrary
coefficient vector, and \cref{step:eqshape} converts it to the separable
$\eq$-shaped vector required by the tensorized F2Z opening.  The interaction is
one row challenge followed by a degree-$2$ sumcheck over the witness-position
variables.

\subsection{Normalize the statement to one witness vector}

The public SHA data must still enter the relation, but it need not be a witness
block.  Let $A_{\xx}$ denote the constraint matrix after substituting the public
message, chaining value, digest, and round constants.  Add one verifier-known
constant coordinate to the source vector,
\[
  \ff^{\circ}=(1,\ff),
\]
and let $M^{\circ}$ contain the constant row, the identity rows for $\ff$, and
the seven XOR families.  Define
\[
  \hh^{\circ}
  =\pitwo^{-1}\!\bigl(M^{\circ}\pitwo(\ff^{\circ})\bigr),
  \qquad
  A_{\xx}^{\circ}\hh^{\circ}=0\quad\text{over }\Z .
\]
The constant column of $A_{\xx}^{\circ}$ carries every affine term created by
substituting $\xx$.  Thus the PIOP witness really is just
\[
  \boxed{\ z=\hh,\qquad A\hh=0,\qquad
  \pitwo(\hh)=M\pitwo(\ff^{\circ})\ },                 \tag{3.1}
\]
where from now on the circles and the subscript $\xx$ are suppressed.  Only the
tail $\ff$ is committed.  The leading $1$ is verifier-known and will reappear as
a known scalar when $M$ is transposed in \cref{step:transpose}.  This
normalization removes the old $(\xx,\ff,\hh)$ concatenation and, with it, the
separate public-input-fold step; it does \emph{not} remove the public input from
the statement.

The two obligations not proved by the zero-check are already fixed.  Commit
time fixed
\[
  C:\quad \orc=\cw{\ff},
\]
and the public virtual map fixes
\[
  H:\quad \pitwo(\hh)=M\pitwo(\ff^{\circ}).
\]
They remain open throughout this section.  The opening IOP will discharge $H$
by the full adjoint of $M$, and $C$ by the final proximity proof.

\paragraph{Indices used below.}
Constraint rows are indexed by
$\beta\in\bits^{\nu}$.  The symbol $\beta$ is deliberately kept separate from
$\rowidx$ and $\colidx$: those two names are reserved for the matrix reshape of
the witness in \cref{sec:iopp}.  Every inner product is annotated by its
ambient field; no inner-product subscript denotes an axis.

\clgreset
\renewcommand{\clgboardtitle}{Open claims}
\clgdef{$\cw{\cdot}$}{$\cw{\ff}=\Enc_{\cC}(\pack(\pitwo(\ff)))$,
  \cref{eq:cw}}

\begin{step*}{Where the PIOP starts}{%
  \clgnew{S}{$A\hh=0$ over $\Z$}%
  \clgnew{C}{$\orc=\cw{\ff}$}%
  \clgnew{H}{$\pitwo(\hh)=M\pitwo(\ff^{\circ})$}%
  \clgnote{S}{$\hh$ is the only PIOP witness vector}%
}
The statement claim $S$ is active.  Claims $C$ and $H$ are standing
obligations: the PIOP carries them but cannot inspect the oracle or prove that a
materialized virtual row came from $M$.
\end{step*}

% =========================================================================
\begin{step}{Multilinearize the constraint rows}{%
  \clgupd{S}{$Q(\beta)=0$ for every $\beta\in\bits^{\nu}$}%
  \clgnote{S}{identity over $\Z$; no challenge and no message}%
}\label{step:mlin}
Pad the $m$ live constraint rows to $2^{\nu}$ and write $A[\beta,*]$ for one
row.  Define the multilinear extension of the residual table:
\[
  Q(Z)
  \defeq
  \sum_{\beta\in\bits^{\nu}}
     \eq(\beta,Z)\,\ipZ{A[\beta,*]}{\hh}.
\]
At a Boolean point, $\eq$ is the Kronecker delta, hence
\[
  A\hh=0
  \quad\Longleftrightarrow\quad
  Q(\beta)=0\quad\text{for every }\beta\in\bits^{\nu}.
\]
This is only a change of notation.  In particular, there is nothing to sum out
yet: the universally quantified row label $\beta$ has merely become the input
of one multilinear polynomial.

\end{step}

% =========================================================================
\begin{step}{Project the integer equations into $\Fq$}{%
  \clgupd{S}{$\bar A\,\pi_q(\hh)=0$ over $\Fq$}%
  \clgnote{S}{$\bar A=\pi_q(A)$; $q=2^{100}-15$ is fixed}%
}\label{step:project}
Reduce the matrix entrywise:
\[
  \bar A\defeq\pi_q(A),
  \qquad
  Q_q(Z)
  \defeq
  \sum_{\beta\in\bits^{\nu}}
    \eq(\beta,Z)\,\ipq{\bar A[\beta,*]}{\pi_q(\hh)}.
\]
All entries of $\hh$ are bits and the nonzero entries of $A$ are signed powers
of two no larger than $2^{35}$.  The public row-wise residual bound is far below
$q$, so reduction cannot turn a nonzero reachable integer residual into zero.

\begin{remark}[Why a fixed $q$ is sound here]\label{warn:fixedq}
The required side condition is
\[
  q>2\max_{\beta}\sum_{\mathsf{pos}}|A[\beta,\mathsf{pos}]|.
\]
For this SHA-256 matrix the right-hand side is below approximately $2^{38}$,
whereas $q\approx2^{100}$.  Therefore projection is injective on every possible
bit-witness residual and introduces zero soundness error.  A redesign of the
arithmetization must recompute this bound.
\end{remark}

\end{step}

% =========================================================================
\begin{step}{Collapse every constraint row at one random point}{%
  \clgdone{S}%
  \clgnew{$E_0$}{$\ipq{\vv}{\pi_q(\hh)}=0$}%
  \clgnote{$E_0$}{$\vv=\bar A^{\!\top}\eq(\cdot,r_{\rm con})$, public}%
}\label{step:collapse}
The verifier samples one point
$r_{\mathrm{con}}\leftarrow\Fq^{\nu}$ and computes
\[
  \vv
  \defeq
  \bar A^{\!\top}\eq(\cdot,r_{\mathrm{con}}),
  \qquad
  \vv[\mathsf{pos}]
  =\sum_{\beta\in\bits^{\nu}}
      \eq(\beta,r_{\mathrm{con}})\bar A[\beta,\mathsf{pos}].
\]

\paragraph{Why evaluation away from the Boolean cube is still zero.}
Let
\[
  e_{\beta}
  \defeq
  \ipq{\bar A[\beta,*]}{\pi_q(\hh)}.
\]
When the constraint claim holds, $e_{\beta}=0$ for every Boolean row label
$\beta$.  Therefore the defining interpolation formula gives
\[
  Q_q(Z)
  =\sum_{\beta\in\bits^{\nu}}\eq(\beta,Z)e_{\beta}
  =\sum_{\beta\in\bits^{\nu}}\eq(\beta,Z)\cdot 0
  =0
  \qquad\text{in }\Fq[Z_1,\ldots,Z_{\nu}].
\]
Thus $Q_q$ is the zero polynomial, not merely a polynomial that happens to
vanish at the Boolean points, and consequently
$Q_q(r_{\mathrm{con}})=0$ for \emph{every}
$r_{\mathrm{con}}\in\Fq^{\nu}$.  Multilinearity is load-bearing here: a general
polynomial such as $Z(Z-1)$ can vanish on the Boolean cube without vanishing
away from it, but $Q_q$ is by construction the unique multilinear extension of
the residual table.

For any witness, whether valid or not, reassociating the sums gives
\[
  Q_q(r_{\mathrm{con}})
  =\sum_{\beta}\eq(\beta,r_{\mathrm{con}})
       \ipq{\bar A[\beta,*]}{\pi_q(\hh)}
  =\ipq{\vv}{\pi_q(\hh)}.
\]
The verifier replaces the full row claim by the randomized scalar claim
\[
  \boxed{\ipq{\vv}{\pi_q(\hh)}=0}.             \tag{$E_0$}
\]
The prover sends nothing.  If a nonzero multilinear residual table is claimed
to be zero, the random evaluation accepts with probability at most
$\nu/|\Fq|$.

\paragraph{What the row collapse does --- and does not do.}
There is no reason to sumcheck the row label $\beta$: $A$ is public and the
witness occurs linearly, so the complete $\beta$-sum moves onto $A$ and becomes
the public vector $\vv$.  But this does \emph{not} make $\vv$ a tensor product
across the later row/column layout.  The next step sumchecks the
\emph{witness-position} variables; confusing those two index sets is exactly
what allowed the missing reduction to go unnoticed.

For $N$ copies of the same circuit, the circuit-dependent portion of this pass
factors through the one-compression matrix.  Public-data-dependent coefficients
still have to be read; the h-only normalization relocates that work into the
construction of $A_{\xx}$ rather than pretending public I/O is free.

\end{step}

% =========================================================================
\begin{step}{Sumcheck the inner product into an evaluation-shaped claim}{%
  \clgdone{$E_0$}%
  \clgnew{$E_{\rm eq}$}{$c\,\mle{h_q}(r_h)=\tau$}%
  \clgnote{$E_{\rm eq}$}{weights are $c\,\eq(\cdot,r_h)$, hence separable}%
}\label{step:eqshape}
Pad $\hh$ and $\vv$ with zeros to a common length $2^{d_h}$, write
$h_q\defeq\pi_q(\hh)$, and let $\mle{\vv}$ and $\mle{h_q}$ denote their
multilinear extensions.  The claim from \cref{step:collapse} is
\[
  0
  =\sum_{p\in\bits^{d_h}}\vv[p]h_q[p]
  =\sum_{p\in\bits^{d_h}}\mle{\vv}(p)\mle{h_q}(p).    \tag{3.2}\label{eq:eqshape-start}
\]
Run the standard degree-$2$ sumcheck on \eqref{eq:eqshape-start}.  At round
$i$, after challenges $r_1,\ldots,r_{i-1}$, the prover sends
\[
\begin{aligned}
  g_i(T)
  ={}&\!\!\sum_{x_{i+1},\ldots,x_{d_h}\in\bits}
    \mle{\vv}(r_1,\ldots,r_{i-1},T,x_{i+1},\ldots,x_{d_h})\\[-2pt]
   &\hspace{2.2cm}\cdot
    \mle{h_q}(r_1,\ldots,r_{i-1},T,x_{i+1},\ldots,x_{d_h}).
\end{aligned}
\]
The verifier checks
\[
  g_i(0)+g_i(1)=s_{i-1},
\]
samples $r_i\leftarrow\Fq$, and sets $s_i=g_i(r_i)$, starting from
$s_0=0$.  After $d_h$ rounds, put
\[
  r_h=(r_1,\ldots,r_{d_h}),
  \qquad \tau=s_{d_h},
  \qquad c=\mle{\vv}(r_h),
\]
where $c$ is public because $\vv$ is public.  The terminal relation is
\[
  \tau=c\,\mle{h_q}(r_h)
      =\ipq{c\,\eq(\cdot,r_h)}{\pi_q(\hh)}.
\]
Thus the sumcheck replaces the arbitrary coefficient vector by the rank-one
evaluation vector:
\[
  \boxed{
  E_{\rm eq}:\quad
  \ipq{c\,\eq(\cdot,r_h)}{\pi_q(\hh)}=\tau.}         \tag{3.3}\label{eq:eqshape-end}
\]
Keeping the scalar $c$ in the coefficient vector avoids division and remains
valid even when $c=0$.  Equivalently, one may fuse
\cref{step:collapse,step:eqshape} into a Spartan-style matrix sumcheck; the
contracted-$\vv$ presentation is the minimal change because the verifier
already performs the sparse public matrix--vector pass.

\end{step}

\subsection{PIOP handoff and error budget}

The SHA-256 PIOP has exactly four moves:
\[
  \begin{aligned}
  A\hh=0\text{ over }\Z
  &\xRightarrow{\text{multilinearize}}
    Q|_{\bits^{\nu}}=0
   \xRightarrow{\pi_q}
    \bar A\pi_q(\hh)=0,\\
  &\xRightarrow{r_{\mathrm{con}}}
    E_0
   \xRightarrow{\text{degree-2 sumcheck}}
    E_{\rm eq}.
  \end{aligned}
\]
Projection contributes no error under \cref{warn:fixedq}.  Evaluating the
multilinear residual contributes at most $\nu/|\Fq|$, and the degree-$2$
sumcheck contributes at most $2d_h/|\Fq|$.  No ideal check,
polynomial-basis evaluation, public-input fold, or bit-decomposition remains.
